%! TEX program = lualatex
\documentclass[a4paper, hidelinks, 11pt]{article}

%%%%%%%%%%%%%%%%% LOAD PACKAGES %%%%%%%%%%%%%%%%

\usepackage{graphicx} % Required for including pictures

\usepackage[british]{babel}
\usepackage{microtype}

% font settings
\usepackage{fontspec}
\usepackage{mathtools}
\usepackage{textcomp}
\usepackage{eurosym}
\usepackage{unicode-math}
\usepackage{enumitem}

\usepackage{soul} % highlighting & strikeout
\usepackage[colorlinks=true,linkcolor=black,anchorcolor=black,citecolor=black,filecolor=black,menucolor=black,runcolor=black,urlcolor=black]{hyperref} % urls

% Page setup
\usepackage[bottom]{footmisc} % force footnote to the bottom of the page
\usepackage{fancyhdr}
\usepackage{lastpage}
\pagestyle{fancy}
\fancyhf{}
\renewcommand{\headrulewidth}{0pt} % remove top rule
\rfoot{\hspace*{0.97\textwidth}\footnotesize \thepage/\pageref{LastPage}}

% \pagestyle{empty} % Suppress headers and footers

\setlength\parindent{0.5cm} % Paragraph indentation

% prevent widows and orphans
\usepackage[defaultlines=2,all]{nowidow}
\usepackage{needspace}

% headings
\usepackage{titlesec}

\titleformat*{\section}{\normalsize\bfseries}
\titleformat*{\subsection}{\normalsize\bfseries}
\titleformat*{\subsubsection}{\normalsize\bfseries}
\titleformat*{\paragraph}{\normalsize\bfseries}
\titleformat*{\subparagraph}{\normalsize\bfseries}

\defaultfontfeatures{Ligatures=TeX}

\setmainfont[
    Extension=.otf,
    UprightFont= *-regular,
    BoldFont=*-bold,
    ItalicFont=*-italic,
    BoldItalicFont=*-bolditalic,
    ]{TeX Gyre Pagella}

\setsansfont[
        Extension=.otf,
        UprightFont= *-regular,
        BoldFont=*-bold,
        ItalicFont=*-italic,
        BoldItalicFont=*-bolditalic,
        ]{TeX Gyre Heros}

\setmathfont{TeX Gyre Pagella Math}

\setmonofont{TeX Gyre Cursor}[Scale=1]


%%%%%%%%%%%%%%%%% BODY %%%%%%%%%%%%%%%%

\begin{document}

\section*{Dear editors,}

We are extremely grateful to the editors and to our two anonymous reviewers for their very positive appraisal of our manuscript entitled \textbf{``Vaccine-induced time- and age-dependent mucosal immunity to gastrointestinal parasite infection''}, who found our study \emph{``[\ldots] of great interest, [\ldots] well conducted and the manuscript is well written''} (Rev \#1), and \emph{``a great example of how to approach a meaningful analysis of mucosal immunity''} (Rev \#2), who \emph{``thoroughly enjoyed the read''}.

\noindent We address each of their suggestions below:


\subsection*{Reviewer \#1}

\begin{itemize}
        \item Line 32: vaccination is not the only promising strategy but rather ``one'' of the.
        \item[\rightarrow] \textcolor{blue}{Modified as suggested.}
        \item Line 361: Please indicate the number of male and female in each group and their status regarding \emph{T. circumcincta} infection prior vaccine immunisation (naive or not?)
        \item[\rightarrow] \textcolor{blue}{All lambs uninfected prior to immunisation. Methods section updated accordingly.}
        \item Line 397: Please indicate the threshold chosen for the RIN (as extraction of good quality RNA from gastrointestinal mucosa is not an easy task)
        \item[\rightarrow] \textcolor{blue}{Indeed, this step was conducted with great care by Mairi Mitchell. The lowest RIN value upon submission to the sequencing facility exceeded 8.2. We added this information to the methods as well as the sequencing Phred scores.}
        \item Finally, to my knowledge, RNAseq data used for expression studies are confirmed by quantitative RT-PCR of random chosen genes (n=8 to 12 generally) followed by estimation of the correlation. If it has been done, please indicate it.
        \item[\rightarrow] \textcolor{blue}{This was not done. However, because each sample was a high quality pool of three biopsies, each lamb was sampled multiple times across time, and our analysis focussed on a pathway analysis rather than on variation in individual gene expression, we felt confident that technical variation would be averaged out. In addition, RNAseq is now a robust technology; when quality checks of input RNA and sequencing quality are good (as they were for our samples), there is little reason to expect significant bias from RNAseq.}
\end{itemize}

% \pagebreak
\subsection*{Reviewer \#2}

\begin{itemize}
        \item 78++ please list actual worm counts and cumulative FWEC and \% protection, this helps gauge the start point for the analysis.
        \item[\rightarrow] \textcolor{blue}{We now include median counts in the main text and have duplicated the \% protection that was provided in the legend of Fig. 1.}
        \item Interesting that an adjuvant control was used as Quil A has substantial inflammatory activity (that worms don't like) and so an unvaccinated control group may have provided a greater contrast.
        \item[\rightarrow] \textcolor{blue}{We fully agree that this would have been informative; indeed, the contrast could be greater between ``no Quil A'' and vaccinated animals. We now acknowledge this and cite Haçariz \emph{et al.}\footnotemark in the discussion. We had decided that the use of 32 additional lambs and corresponding transcriptome costs would not be justified by a relatively secondary question, given that we were mostly interested in the antigen-specific effects of the prototype vaccine. We now clarify this in the methods section.}
        \item 95- please clarify that the overall analyses were taken to establish protective effects to reduce the worm count or cumulative egg count irrespective of group, with vaccination/ age as internal variables. I state this as the thrust of the paper is to identify those protective pathways to both vaccine success but also parasite resistance through innate and acquired immunity.
        \item[\rightarrow] \textcolor{blue}{Many thanks for this suggestion -- it greatly clarifies the presentation of the results.}
        \item Wow, 5mg QuilA is quite a high dose!!
        \item[\rightarrow] \textcolor{blue}{We realise this is a very high dose and have reduced the dose to 1mg for subsequent vaccine studies.}
        \item 207. What effect can be attributed to the QuilA vs the vaccine proteins at day0. Std controls without QuilA might have given a different starting point.
        \item[\rightarrow] \textcolor{blue}{As mentioned above, we cannot estimate the QuilA-only effect on worm burdens in this experiment. However, it could have either no effect, an effect that depends on the age group, or that is independent of age. Given the age-independence of innate responses, we would hypothesise that any QuilA-specific effects should be constant across age groups. However, this is something that future studies will need to test. We have added a sentence accordingly in the discussion.}
        \item 213--233: I contend that this section could be omitted as simply repeats earlier sections. Real discussion commences around line 251
        \item[\rightarrow] \textcolor{blue}{We agree, and have cut text that was too repetitive in first paragraph of the discussion.}
        \item 278+ could the increased levels of innate and inflammatory responses be related to ongoing worm damage to the mucosa as the cumulative life cycles are effected?
        \item[\rightarrow] \textcolor{blue}{We thank the reviewer for this important insight. Indeed, this seems likely. We have added this hypothesis to the discussion accordingly.}
\end{itemize}

% \closing{Sincerely, on behalf of my co-authors,}

% FOOTNOTES
\footnotetext{Hacariz, O. et al. The effect of Quil A adjuvant on the course of experimental Fasciola hepatica infection in sheep. Vaccine 27, 45-50 (2009).}

\end{document}
